\documentclass[11pt]{article}

\usepackage[margin=1in]{geometry}
\usepackage{amsmath,amssymb}
\usepackage{booktabs}
\usepackage{longtable}
\usepackage{array}
\usepackage{enumitem}
\usepackage{xcolor}
\usepackage[strings]{underscore}
\usepackage[hidelinks]{hyperref}

\setlist[itemize]{leftmargin=*}
\setlist[enumerate]{leftmargin=*}
\newcommand{\OffCEM}{\textsc{OffCEM}}
\newcommand{\IPS}{\textsc{IPS}}
\newcommand{\MIPS}{\textsc{MIPS}}
\newcommand{\DM}{\textsc{DM}}
\newcommand{\DR}{\textsc{DR}}
\newcommand{\SNIPS}{\textsc{SNIPS}}
\newcommand{\relMSE}{\mathrm{relMSE}}
\newcommand{\relBias}{\mathrm{relBias}^2}
\newcommand{\relVar}{\mathrm{relVar}}
\newcommand{\dm}{\mathrm{dm}}

\title{OffCEM Experimental Setup and Completed Experiments}
\author{OffCEM extension workspace}
\date{Compiled from workspace notes, 2026-06-17}

\begin{document}
\maketitle

\section{Scope and Sources}

This document summarizes the experimental setup and the experiments with results
currently recorded in the OffCEM extension workspace. It is based on the local
markdown notes and result write-ups, especially:

\begin{itemize}
  \item \texttt{icml2023-offcem-expts/handoff.md}
  \item \texttt{icml2023-offcem-expts/experimental_setup.md}
  \item \texttt{icml2023-offcem-expts/README.md}
  \item \texttt{local_correctness_results/local_correctness_findings.md}
  \item \texttt{clustering_results/clustering_experiment_findings.md}
  \item \texttt{imprecise_results/imprecise_findings.md}
  \item \texttt{icml2023-offcem-expts/src/synthetic/applicability/README.md}
  \item \texttt{icml2023-offcem-expts/src/synthetic/imprecise/README.md}
  \item \texttt{imprecise_clustering_methods.md}
  \item \texttt{offcem_paper_breakdown.md}
\end{itemize}

Prepared-only manifests and planned runs are intentionally not listed as
completed experiments. When older results are known to have been produced by
superseded code, this is stated explicitly.

\section{Experimental Setup}

\subsection{Research Question}

The base paper studies off-policy evaluation (OPE) for contextual bandits with a
large discrete action set. A logging policy $\pi_0$ generated a dataset
$\mathcal{D} = \{(x_i,a_i,r_i)\}_{i=1}^n$, and the goal is to estimate the value
of a target policy $\pi$,
\[
  V(\pi) = \mathbb{E}_{p(x)\pi(a\mid x)}[q(x,a)],
  \qquad
  q(x,a) = \mathbb{E}[r\mid x,a].
\]
The OffCEM extension work asks how the action clustering used by \OffCEM{}
affects bias, variance, local correctness, and whether OffCEM beats standard OPE
baselines.

\subsection{Software Environment}

Experiments are run in the \texttt{offcem} conda environment with Python 3.9 and
OBP 0.5.5. The base environment is not suitable because it has an incompatible
OBP version. Synthetic modules are run from:
\[
  \texttt{icml2023-offcem-expts/src/synthetic}.
\]
The preferred command shape is:
\[
  \texttt{conda run -n offcem python -m <module> ...}
\]
Long experiment sweeps are controlled manually; the notes permit quick smoke
tests, syntax checks, focused unit tests, and analysis-only commands.

\subsection{Default Synthetic Bandit World}

The default benchmark is a fully synthetic contextual bandit environment. It is
used because the true policy value can be computed and compared against OPE
estimates.

\begin{table}[h]
\centering
\begin{tabular}{ll}
\toprule
Parameter & Default value \\
\midrule
Number of logged rounds $n$ & 3000 \\
Training users in original synthetic setup & 200 \\
Ground-truth test rounds & 100000 \\
Number of actions & 1000 \\
Context dimension & 10 \\
Action embedding dimension & 10 categorical dimensions \\
Action embedding cardinality & 5 values per categorical dimension \\
Number of observed clusters & 50 by default \\
Number of unobserved categorical dimensions & 0 \\
Reward noise standard deviation & 3.0 \\
Logging-policy parameter $\beta$ & -0.1 \\
Target-policy exploration $\epsilon$ & 0.2 \\
Deficient actions & 0 in the default setting \\
Test users in code & 1000 \\
Default synthetic random state & 12345 \\
\bottomrule
\end{tabular}
\caption{Default synthetic setup recorded in the workspace notes. The notes flag
one paper/code discrepancy: \texttt{N\_TEST\_USERS}=1000 in code versus 200 in
the paper. This affects ground-truth computation only, not logged OPE data.}
\end{table}

Each context $x\in\mathbb{R}^{10}$ is sampled from a standard normal
distribution. Each action has a synthetic categorical embedding: ten categorical
attributes with five possible values each. The one-hot expansion of these
features is the representation used by most clustering methods and reward
models.

The paper-breakdown notes record two implementation details that matter for
reproducibility. First, OBP 0.5.5 internally increments
\texttt{n\_cat\_dim} and \texttt{n\_unobserved\_cat\_dim} by one, so a configured
10-dimensional observed categorical embedding is represented internally with an
additional unobserved categorical dimension. Second, action embeddings are
stochastic: the environment stores a probability matrix $p(e\mid a)$ rather than
only hard action features, and MIPS marginalizes over this distribution. The
sampled embeddings are regenerated when \texttt{obtain\_batch\_bandit\_feedback}
is called, so training and test environments are draws from the same population
rather than the exact same finite environment.

\subsection{Reward Generation}

The expected reward is constructed as a conjunct effect model:
\[
  q(x,a) = g(x,c_a) + h_{c_a}(x,a),
  \qquad c_a = \phi_{\mathrm{gen}}(a).
\]
The reward-generating partition $\phi_{\mathrm{gen}}$ is therefore part of the
data-generating process.

\paragraph{Cluster effect.}
The cluster-level term $g(x,c)$ is constant across actions in cluster $c$ for a
fixed context. It is built from a degree-3 polynomial interaction between the
context and the cluster identity, plus four nonlinear indicator terms in the
context. In the plain-language setup note, these are described as broad
context-cluster preferences.

\paragraph{Residual effect.}
The within-cluster residual $h_c(x,a)$ captures action-specific variation within
the cluster:
\[
  h_c(x,a)
  =
  x^\top M_c \,\mathrm{onehot}(a)
  + \theta_{x,c}^\top x
  + \theta_{a,c}^\top \mathrm{onehot}(a).
\]
Each cluster receives independent residual parameters using a cluster-indexed
random seed, recorded in the notes as \texttt{random\_state + c}. This detail is
central: in the default benchmark, the residual structure is keyed to the
generating cluster index rather than to a stable feature geometry.

\paragraph{Observed reward.}
The logged reward is:
\[
  r = q(x,a) + \varepsilon,
  \qquad \varepsilon \sim \mathcal{N}(0,3^2).
\]

\subsection{Policies}

The logging policy is a softmax over expected rewards:
\[
  \pi_0(a\mid x)
  =
  \frac{\exp(\beta q(x,a))}
       {\sum_{a'}\exp(\beta q(x,a'))},
  \qquad \beta=-0.1.
\]
Because $\beta$ is negative, the logging policy is intentionally suboptimal and
slightly favors worse actions.

The target policy is $\epsilon$-greedy:
\[
  \pi(a\mid x)
  =
  (1-\epsilon)\mathbf{1}\{a=a^*(x)\}+\frac{\epsilon}{|\mathcal{A}|},
  \qquad \epsilon=0.2,
\]
where $a^*(x)$ is the action with the largest expected reward.

\subsection{OffCEM Estimator}

OffCEM uses a cluster-level importance weight:
\[
  w(x,c)
  =
  \frac{\pi(c\mid x)}{\pi_0(c\mid x)}
  =
  \frac{\sum_{a:\phi(a)=c}\pi(a\mid x)}
       {\sum_{a:\phi(a)=c}\pi_0(a\mid x)}.
\]
The estimator has the same control-variate shape as doubly robust OPE, but with
cluster weights:
\[
  \hat V_{\mathrm{OffCEM}}
  =
  \frac{1}{n}\sum_{i=1}^n
  \left[
    w(x_i,\phi(a_i))(r_i-\hat f(x_i,a_i))
    + \hat f(x_i,\pi)
  \right],
\]
where
\[
  \hat f(x,\pi)=\sum_a \pi(a\mid x)\hat f(x,a).
\]

\subsection{Two-Stage Reward Regression}

The two-stage OffCEM reward model decomposes:
\[
  \hat f(x,a)=\hat g_\psi(x,\phi(a))+\hat h_\theta(x,a).
\]
The first stage fits within-cluster reward differences using action pairs in the
same cluster:
\[
  \min_\theta
  \sum_{(x,a,b)}
  \ell_h\left(r_a-r_b,\,
  \hat h_\theta(x,a)-\hat h_\theta(x,b)\right).
\]
This directly targets within-cluster contrast accuracy. The second stage freezes
$\hat h_\theta$ and fits a cluster-level offset:
\[
  \min_\psi
  \sum_{(x,a,r)}
  \ell_g\left(r-\hat h_\theta(x,a),\,
  \hat g_\psi(x,\phi(a))\right).
\]
The second stage cannot change within-cluster contrasts because the cluster
offset cancels for any two actions in the same cluster.

\subsection{Reward-Model Implementation Details}

The paper-breakdown note records the following concrete model choices from the
code.

\paragraph{Pairwise model.}
The first-stage pairwise model is a PyTorch \texttt{PairWiseRegression} network.
It maps a 10-dimensional context vector to a 1000-dimensional vector of
per-action predictions. The synthetic architecture is four ELU layers with
dimensions:
\[
  10 \rightarrow 30 \rightarrow 30 \rightarrow 30 \rightarrow 1000.
\]
Actions are represented by output index rather than by feeding action features
into the network. Training uses within-cluster action pairs sharing the same user,
AdamW with learning rate 0.01 and weight decay $10^{-4}$, batch size 128, 30
epochs, and an exponential learning-rate scheduler with factor 0.95.

\paragraph{Cluster-offset model.}
The second-stage cluster offset uses OBP's \texttt{RegressionModel} wrapping an
sklearn \texttt{MLPRegressor} with hidden layers $(50,50,50)$. Its input is the
concatenation of context and one-hot cluster identity, and its target is
$r-\hat h_\theta(x,a)$.

\paragraph{Baseline reward model.}
The direct-method and doubly-robust reward model uses the same
\texttt{MLPRegressor(50,50,50)} pattern with action context as one-hot action
features. The code uses \texttt{n\_folds=1}, so there is no cross-fitting in the
implementation despite cross-fitting being discussed in the paper.

\paragraph{Result accounting.}
The code stores per-seed estimates, then computes bias, variance, MSE, and
normalized versions downstream. Older paper-reproduction code stores columns
named ``bias'' and ``variance'' as per-seed components that are only meaningful
after aggregation.

\subsection{Local Correctness}

OffCEM is unbiased under common cluster support and local correctness. Local
correctness requires the reward model to preserve true within-cluster reward
differences:
\[
  q(x,a)-q(x,b)=\hat f(x,a)-\hat f(x,b)
  \quad\text{for all }a,b\text{ with }\phi(a)=\phi(b).
\]
Equivalently, the regression error $q(x,a)-\hat f(x,a)$ may be wrong in absolute
level, but it must be constant within each cluster.

The local-correctness experiments use a demeaned error metric. For model error
$e(x,a)=\hat f(x,a)-q(x,a)$,
\[
  \bar e_\phi(x,c)=\frac{1}{|A_c|}\sum_{a'\in A_c} e(x,a'),
  \qquad
  e_{\mathrm{dm}}(x,a)=e(x,a)-\bar e_\phi(x,\phi(a)),
\]
and
\[
  \dm(\hat f;\phi)
  =
  \mathbb{E}_x\left[
  \frac{1}{|\mathcal{A}|}\sum_a e_{\mathrm{dm}}(x,a)^2
  \right].
\]
The ratio $\dm_{2s}/\dm_{1s}$ compares two-stage OffCEM against a one-stage
reward model. Values below one mean two-stage training improves within-cluster
relative reward prediction.

\subsection{Clustering Methods}

The workspace implements and evaluates several hard partitions. Unless stated
otherwise, each method receives the one-hot action-feature matrix, the requested
cluster count $K$, and the experiment seed. After the base assignment, the shared
\texttt{compute\_clusters} entrypoint dense-relabels cluster IDs into consecutive
values. In the clustering experiment, base assignments can also be post-processed
with a balance mode: \texttt{natural} leaves method output unchanged;
\texttt{balanced} assigns target sizes as evenly as possible; and
\texttt{power\_law} uses Zipf weights with exponent $\alpha=1.0$. Balance
post-processing sorts actions by the original cluster labels and then slices that
stable ordering into the requested target sizes.

\begin{itemize}
  \item \texttt{original}: the paper/code method, implemented as a random linear
  score over action features, softmax with temperature, and stochastic sampling.
  The default temperature is 10. The notes emphasize that this is weakly
  feature-dependent but not geometrically coherent.

  \item \texttt{kmeans}: sklearn \texttt{KMeans} on the one-hot action-feature
  matrix, called as \texttt{KMeans(n\_clusters=K, random\_state=seed,
  n\_init=10)} under scikit-learn 1.0.2. No kernel is used. The objective is
  ordinary Euclidean squared distance to cluster centroids. Other sklearn
  defaults apply: \texttt{init="k-means++"}, \texttt{max\_iter=300},
  \texttt{tol=1e-4}, \texttt{copy\_x=True}, and \texttt{algorithm="auto"}.

  \item \texttt{spectral}: sklearn \texttt{SpectralClustering} with
  \texttt{affinity="nearest\_neighbors"}, \texttt{n\_neighbors=min(20,
  n\_actions-1)}, \texttt{assign\_labels="kmeans"}, and the experiment
  \texttt{random\_state}. This constructs a nearest-neighbor graph from the
  action-feature matrix, embeds via the graph Laplacian, and labels the spectral
  embedding with K-means. Other sklearn defaults apply, including
  \texttt{eigen\_solver=None}, \texttt{n\_components=None},
  \texttt{n\_init=10}, \texttt{eigen\_tol=0.0}, and \texttt{n\_jobs=None}.

  \item \texttt{agglomerative}: sklearn \texttt{AgglomerativeClustering} with
  \texttt{n\_clusters=K} and \texttt{linkage="ward"}. Ward linkage uses a
  Euclidean variance-minimization criterion. No kernel is supplied. Other
  sklearn defaults apply: \texttt{affinity="euclidean"},
  \texttt{connectivity=None}, \texttt{compute\_full\_tree="auto"},
  \texttt{distance\_threshold=None}, and \texttt{compute\_distances=False}.

  \item \texttt{random}: independent random cluster labels sampled with
  \texttt{check\_random\_state(seed).randint(0,K,size=n\_actions)}. Cluster sizes
  are therefore multinomial and may be imbalanced before any optional balance
  post-processing.

  \item \texttt{shuffled\_partition}: random equal-size partition. The method
  permutes action indices with \texttt{check\_random\_state(seed).permutation},
  then slices the permutation into $K$ contiguous groups. Each group has size
  $\lfloor n_{\mathrm{actions}}/K\rfloor$ or one extra action for the first
  remainder groups. This is the size-balanced random control.

  \item \texttt{feature\_bucket}: deterministic feature-only clustering by
  mean-centering the action-feature matrix, computing an SVD, projecting actions
  onto the first right singular vector, ranking projected values with ordinal
  ranks, and splitting the ranking into $K$ equal-frequency bins via
  $\lfloor \mathrm{rank}(a)K/n_{\mathrm{actions}}\rfloor$. The \texttt{random\_state}
  argument is accepted only for API compatibility; the assignment is deterministic
  given the action features.

  \item \texttt{corrupt\_p}: a controlled corruption of the matched partition by
  selecting $\mathrm{round}(p\,n_{\mathrm{actions}})$ actions without replacement
  and reassigning each selected action to an independent uniform label in
  $\{0,\ldots,K-1\}$. A selected action can receive its original label, so the
  expected changed fraction is $p(1-1/K)$. The output is dense-relabelled after
  corruption.
\end{itemize}

For imprecise-clustering pilots, hard-label adapters and diagnostics also exist
for fuzzy C-means (FCM), possibilistic C-means (PCM), and rough K-means. ECM is
documented as optional until \texttt{evclust} is installed and pinned.

\subsection{Estimators and Metrics}

The standard estimator set across the experiments includes \IPS{}, \DM{}, \DR{},
\MIPS{}, OffCEM one-stage, and OffCEM two-stage. Some newer harnesses also
include cluster-IPS and \SNIPS{}.

For estimator accuracy, the common normalized metrics are:
\[
  \relMSE = \frac{\mathrm{Bias}^2+\mathrm{Var}}{V(\pi)^2},
  \qquad
  \relBias = \frac{\mathrm{Bias}^2}{V(\pi)^2},
  \qquad
  \relVar = \frac{\mathrm{Var}}{V(\pi)^2}.
\]
The local-correctness experiments additionally report $\dm_{2s}$, $\dm_{1s}$,
and $\dm_{2s}/\dm_{1s}$.

\section{Completed Experiments With Results}

\subsection{Clustering Experiment: Matched and Mismatched Modes}

\paragraph{Purpose.}
This experiment studied how clustering method, cluster count, and cluster balance
affect OffCEM OPE performance.

\paragraph{Result paths.}
\begin{itemize}
  \item \texttt{/Users/cispa/Documents/OffCEM/clustering\_results/matched/}
  \item \texttt{/Users/cispa/Documents/OffCEM/clustering\_results/mismatched/}
  \item \texttt{/Users/cispa/Documents/OffCEM/clustering\_results/clustering\_experiment\_findings.md}
\end{itemize}

\paragraph{Design.}
Matched mode used the same clustering method for reward generation and
estimation. Completed matched results cover \texttt{original}, \texttt{kmeans},
\texttt{spectral}, and \texttt{random}; cluster counts 10, 50, and 200; natural
and balanced cluster balance; 15 seeds. Agglomerative was excluded from the main
matched grid because it was too slow, with only partial output.

Mismatched mode generated data with \texttt{original}/natural clustering and
estimated with alternative clusterings. It used cluster counts 10, 50, and 200
with 60 seeds. The notes report all five estimation methods at $n_c=200$, with
some methods absent at lower cluster counts due to the run structure.

\paragraph{Representative baseline comparison.}
In the reference setting \texttt{original}, $n_c=50$, natural balance:
\begin{table}[h]
\centering
\begin{tabular}{lrrrr}
\toprule
Estimator & Mean & Bias & \relMSE & \relVar \\
\midrule
\IPS{} & 3.84 & -0.27 & 1.500 & 1.495 \\
\MIPS{} & 3.79 & -0.32 & 1.480 & 1.473 \\
\DR{} & -0.81 & -4.92 & 1.483 & 0.048 \\
\DM{} & -1.20 & -5.31 & 1.672 & 0.003 \\
OffCEM 1-stage & -1.02 & -5.13 & 1.564 & 0.005 \\
OffCEM 2-stage & 0.92 & -3.19 & 0.605 & 0.001 \\
\bottomrule
\end{tabular}
\end{table}

\paragraph{Findings.}
Cluster count was the dominant factor in matched mode: OffCEM 2-stage relMSE
roughly doubled from $n_c=10$ to $n_c=200$. No learned method consistently beat
the original clustering in the matched results; differences among matched
methods were much smaller than the cluster-count effect. Random clustering
degraded at higher cluster counts. Bias, not variance, dominated OffCEM error.

In mismatched mode, non-original estimation clusterings performed much worse
than the original estimation clustering. At $n_c=10$, OffCEM with \texttt{kmeans}
had relMSE 1.692 versus 0.466 for \texttt{original}. At $n_c=200$, all
non-original methods were clustered around relMSE 1.53--1.60, while
\texttt{original} was 1.019.

\paragraph{Caveat.}
The local-correctness write-up later found two bugs affecting the older
clustering-experiment code: non-reproducible original clustering and computing
$V_{\mathrm{true}}$ at the wrong cluster count. Therefore these results are
listed as completed runs, but their magnitudes should be regenerated before
being treated as final.

\subsection{Local Correctness: Cluster-Count Sweep}

\paragraph{Purpose.}
This experiment directly tested whether the OffCEM two-stage reward model
preserves within-cluster relative reward differences.

\paragraph{Result paths.}
\texttt{/Users/cispa/Documents/OffCEM/local\_correctness\_results/\{matched,original,random,shuffled\_partition\}/nc\{10,50,100,200\}.json}
and the associated PNG diagnostics.

\paragraph{Design.}
The sweep used $n_c\in\{10,50,100,200\}$ and methods
\texttt{matched}, \texttt{original}, \texttt{random}, and
\texttt{shuffled\_partition}. The default setup used 3000 logged rounds, 1000
actions, reward noise 3, target $\epsilon=0.2$, and 10 seeds.

\paragraph{Result.}
\begin{table}[h]
\centering
\begin{tabular}{lrrrr}
\toprule
Method & $n_c=10$ & $n_c=50$ & $n_c=100$ & $n_c=200$ \\
\midrule
\texttt{matched} & 0.61 & 0.41 & 0.20 & 0.14 \\
\texttt{original} & 0.93 & 0.72 & 0.46 & 0.39 \\
\texttt{random} & 0.97 & 0.75 & 0.46 & 0.38 \\
\texttt{shuffled\_partition} & 0.94 & 0.68 & 0.46 & 0.36 \\
\bottomrule
\end{tabular}
\caption{Local-correctness ratio $\dm_{2s}/\dm_{1s}$. Lower is better.}
\end{table}

\paragraph{Findings.}
The exact matched partition had the strongest local-correctness benefit at every
cluster count. Mismatched partitions retained a weaker two-stage benefit, but the
large matched advantage disappeared. More clusters reduced within-cluster
heterogeneity and lowered the demeaned error, even though separate relMSE
experiments showed that more clusters can worsen OPE bias/variance. Size
imbalance mainly affected tail risk: larger clusters had higher per-cluster
demeaned error, making random partitions worse in their worst clusters.

\subsection{Local Correctness: Temperature Sweep}

\paragraph{Purpose.}
This tested whether lowering the original clustering temperature makes the
generating partition more feature-coherent and therefore recoverable by
\texttt{kmeans}.

\paragraph{Result paths.}
\texttt{/Users/cispa/Documents/OffCEM/local\_correctness\_results/temp\_<T>/<method>/nc50.json}
for $T\in\{1,2,5,10\}$.

\paragraph{Design.}
Cluster count was fixed at 50. Methods were \texttt{matched}, \texttt{kmeans},
\texttt{original}, and \texttt{random}.

\paragraph{Result.}
\begin{table}[h]
\centering
\begin{tabular}{lrrrr}
\toprule
Method & $T=1$ & $T=2$ & $T=5$ & $T=10$ \\
\midrule
\texttt{matched} & 0.46 & 0.41 & 0.46 & 0.44 \\
\texttt{kmeans} & 0.84 & 0.83 & 0.78 & 0.79 \\
\texttt{original} & 0.71 & 0.70 & 0.71 & 0.74 \\
\texttt{random} & 0.74 & 0.73 & 0.70 & 0.76 \\
\bottomrule
\end{tabular}
\caption{Local-correctness ratio by original-clustering temperature.}
\end{table}

\paragraph{Findings.}
The hypothesis was refuted. Results were nearly flat across temperature.
\texttt{kmeans} was worse than \texttt{random} at every temperature. Adjusted
Rand index between the generating partition and K-means remained near zero even
at very low temperatures. Temperature changes cluster-size balance, not
feature-geometric coherence.

\subsection{Local Correctness: Corruption Sweep}

\paragraph{Purpose.}
This experiment created a continuous path from the exact matched partition to a
nearly random partition by corrupting a fraction $p$ of action labels.

\paragraph{Result paths.}
\texttt{/Users/cispa/Documents/OffCEM/local\_correctness\_results/corrupt\_p\_<p>/corrupt/nc50.json}.

\paragraph{Design.}
Cluster count was fixed at 50. The corruption levels were
$p\in\{0,0.1,0.25,0.5,0.75,1\}$ with 10 seeds. A selected action can be assigned
its original label, so the expected changed fraction is $p(1-1/K)$.

\paragraph{Result.}
\begin{table}[h]
\centering
\begin{tabular}{lrr}
\toprule
$p$ & $\dm_{2s}/\dm_{1s}$ & Approx. within-cluster error reduction \\
\midrule
0.00 & 0.437 & 56\% \\
0.10 & 0.491 & 51\% \\
0.25 & 0.580 & 42\% \\
0.50 & 0.633 & 37\% \\
0.75 & 0.684 & 32\% \\
1.00 & 0.713 & 29\% \\
\bottomrule
\end{tabular}
\end{table}

\paragraph{Findings.}
Local-correctness violation degraded smoothly and monotonically; there was no
binary cliff. The two-stage OffCEM point estimate also degraded monotonically as
the matched partition was corrupted. Standard baselines were flat across $p$
because they do not use the clustering.

\subsection{Feature-Coherent Local Correctness Benchmark}

\paragraph{Purpose.}
This benchmark tested whether the fragility observed on the default
feature-arbitrary reward is intrinsic to OffCEM or an artifact of the synthetic
reward being keyed to a non-geometric partition.

\paragraph{Result paths.}
\begin{itemize}
  \item \texttt{/Users/cispa/Documents/OffCEM/local\_correctness\_results/feature\_coherent\_fb/}
  \item \texttt{/Users/cispa/Documents/OffCEM/local\_correctness\_results/feature\_coherent\_km/}
\end{itemize}

\paragraph{Design.}
The reward-generating clustering was changed using
\texttt{--gen-clustering-method}. The primary run generated rewards using the
deterministic \texttt{feature\_bucket} partition and estimated with
\texttt{matched}, \texttt{feature\_bucket}, \texttt{kmeans}, and \texttt{random}.
A contrast run generated rewards using \texttt{kmeans}. Both used $n_c=50$ and
10 seeds.

\paragraph{Primary result.}
\begin{table}[h]
\centering
\begin{tabular}{lrrr}
\toprule
Estimation method & $\dm_{2s}$ & $\dm_{1s}$ & Ratio \\
\midrule
\texttt{matched} & 7.21 & 18.06 & 0.40 \\
\texttt{feature\_bucket} & 7.00 & 17.96 & 0.39 \\
\texttt{kmeans} & 17.51 & 20.74 & 0.85 \\
\texttt{random} & 16.95 & 22.77 & 0.74 \\
\bottomrule
\end{tabular}
\end{table}

\paragraph{Findings.}
When the reward genuinely follows an observable feature geometry,
\texttt{feature\_bucket} recovers oracle-level local correctness without access to
the generating labels. K-means still fails because it captures a different
feature geometry from the one that generated the reward. The refined conclusion
is that local correctness depends on alignment between the estimation partition
and the reward residual geometry; that alignment can be achieved without oracle
labels, but feature use alone is insufficient.

\subsection{Join Experiment: Does Local-Correctness Violation Predict OPE Error?}

\paragraph{Purpose.}
This experiment paired local-correctness violation and OffCEM estimator error on
identical runs. It tested whether $\dm$ is a useful within-problem diagnostic for
choosing a clustering.

\paragraph{Result path.}
\texttt{/Users/cispa/Documents/OffCEM/join\_results/cell\_*.json}.

\paragraph{Design.}
The grid contained:
\[
  n_c\in\{10,50,200\},\quad
  \epsilon\in\{0,0.2,0.7\},\quad
  \sigma\in\{1,3\},\quad
  n\in\{1000,3000\}.
\]
For each of the 36 cells, the experiment evaluated six clustering conditions:
\texttt{matched}, \texttt{kmeans}, \texttt{random},
\texttt{shuffled\_partition}, \texttt{corrupt\_p0.25}, and
\texttt{corrupt\_p0.75}, with 10 seeds.

\paragraph{Findings.}
Within a fixed problem cell, $\dm$ strongly predicted OffCEM error:
\begin{itemize}
  \item mean Spearman correlation of $\dm$ with $\relBias$ was +0.80;
  \item median Spearman correlation was 0.83;
  \item correlation was positive in 36/36 cells;
  \item $\arg\min(\dm)$ matched $\arg\min(\relBias)$ in 36/36 cells.
\end{itemize}
Across different problem cells, raw $\dm$ did not port because it is in
reward-units squared while normalized estimator error also depends on policy,
cluster count, reward scale, and true value.

\paragraph{Baseline crossover.}
The same grid showed that clustering quality decides whether OffCEM is worth
using relative to a baseline such as DR:
\begin{table}[h]
\centering
\begin{tabular}{lrr}
\toprule
OffCEM 2-stage clustering & Pooled relMSE & Cells beating DR \\
\midrule
\texttt{matched} & 1.32 & 36/36 \\
\texttt{corrupt\_p0.25} & 1.78 & 36/36 \\
\texttt{corrupt\_p0.75} & 2.41 & 16/36 \\
\texttt{random}/\texttt{kmeans}/\texttt{shuffled} & 2.50--2.56 & 6--14/36 \\
\bottomrule
\end{tabular}
\end{table}

\subsection{Join Experiment on Feature-Coherent Rewards}

\paragraph{Purpose.}
This repeated the join grid with the reward generated by the deterministic
\texttt{feature\_bucket} partition. It tested whether oracle-free feature-aligned
clustering reaches oracle estimator error at full grid scale.

\paragraph{Result path.}
\texttt{/Users/cispa/Documents/OffCEM/join\_results\_fc/cell\_*.json}.

\paragraph{Design.}
The grid was the same 36-cell design as the original join. The six clustering
conditions were \texttt{matched}, \texttt{feature\_bucket}, \texttt{kmeans},
\texttt{random}, \texttt{corrupt\_p0.25}, and \texttt{corrupt\_p0.75}.

\paragraph{Findings.}
Within-cell diagnostic performance remained strong: mean Spearman correlation of
$\dm$ with $\relBias$ was +0.85, positive in 36/36 cells. The strict argmin match
dropped because \texttt{matched} and \texttt{feature\_bucket} were genuine
co-winners, not because the diagnostic failed.

\begin{table}[h]
\centering
\begin{tabular}{lrr}
\toprule
OffCEM 2-stage clustering & Pooled relMSE & Cells beating DR \\
\midrule
\texttt{matched} & 1.34 & 36/36 \\
\texttt{feature\_bucket} & 1.34 & 36/36 \\
\texttt{corrupt\_p0.25} & 1.80 & 36/36 \\
\texttt{corrupt\_p0.75} & 2.46 & 7/36 \\
\texttt{kmeans} & 2.49 & 8/36 \\
\texttt{random} & 2.57 & 0/36 \\
\bottomrule
\end{tabular}
\end{table}

The feature-bucket clustering tied the oracle matched partition without using
oracle labels. This is the strongest positive result recorded in the notes.

\subsection{Applicability Benchmark Phase A Smoke}

\paragraph{Purpose.}
The applicability package asks when OffCEM wins or loses against baselines across
a broader synthetic design. It separates reward variance into four components:
\[
  q(x,a)=
  \sqrt{\alpha_g}g(x,\mathrm{cluster}(a))
  +\sqrt{\alpha_f}h_{\mathrm{feature}}(x,a)
  +\sqrt{\alpha_i}h_{\mathrm{identity}}(x,a)
  +\sqrt{\alpha_u}h_{\mathrm{hidden}}(x,a).
\]

\paragraph{Result path.}
\texttt{/Users/cispa/Documents/OffCEM/applicability\_results/smoke/}.

\paragraph{Design.}
Phase A smoke used 2 Sobol cells, 5 seeds, and 6 partitions. The dataset was
generated once per cell/seed and reused across partitions and estimators. The
estimator harness included OffCEM 2-stage, OffCEM 1-stage, cluster-IPS, MIPS,
DR, DM, IPS, and SNIPS. Diagnostics included DM variants, effective sample size,
importance-weight tails, unsupported target mass, partition ARI, cluster-size
statistics, observed actions and pairs per cluster, and realized reward variance
shares.

\paragraph{Findings and caveat.}
The smoke run completed end-to-end and wrote checkpoints,
\texttt{estimator\_metrics.csv}, and \texttt{paired\_comparisons.csv}. The
handoff notes report that OffCEM 2-stage underperformed baselines in this tiny
configuration, even on the matched partition. This is treated as a likely
degenerate-smoke artifact rather than a final research conclusion.

\subsection{Membership-Uncertainty Direction 1: Hardened Clustering Diagnostics}

\paragraph{Purpose.}
This pilot asked whether uncertainty-aware clustering methods improve OffCEM or
diagnose OffCEM failure when hardened to ordinary cluster labels.

\paragraph{Result path.}
\texttt{/Users/cispa/Documents/OffCEM/imprecise\_results/direction1/}.

\paragraph{Design.}
The completed pilot used two generating setups
(\texttt{original}, \texttt{feature\_bucket}), 5 seeds, 1000 actions, 50
clusters, 3000 rounds, and methods \texttt{matched}, \texttt{original},
\texttt{feature\_bucket}, \texttt{kmeans}, \texttt{random}, FCM, PCM, and rough
K-means.

\paragraph{Findings.}
\begin{table}[h]
\centering
\begin{tabular}{lrrrrrr}
\toprule
Generating setup & matched & kmeans & rough & fcm & pcm & random \\
\midrule
\texttt{feature\_bucket} & 0.581 & 0.942 & 0.947 & 1.046 & 1.300 & 1.336 \\
\texttt{original} & 0.770 & 0.951 & 1.114 & 1.172 & 1.355 & 1.133 \\
\bottomrule
\end{tabular}
\caption{OffCEM relMSE from the membership-uncertainty diagnostics pilot.}
\end{table}

The matched partition was the best OffCEM variant in both setups, but standard
baselines were better in the original generating setup. FCM produced nearly
uniform memberships, PCM collapsed all centers, and rough K-means marked about
75\% of actions as boundary objects. Policy-weighted DM ranked OffCEM error
better than raw DM.

\subsection{Membership-Uncertainty Direction 6: Partition Sensitivity}

\paragraph{Purpose.}
This experiment sampled plausible hard partitions from FCM, PCM, and rough
K-means uncertainty outputs to measure sensitivity of OffCEM estimates to the
partition draw.

\paragraph{Result paths.}
\begin{itemize}
  \item \texttt{/Users/cispa/Documents/OffCEM/imprecise\_results/direction6/}
  \item \texttt{/Users/cispa/Documents/OffCEM/imprecise\_results/direction6\_upd/}
\end{itemize}

\paragraph{Design.}
The corrected run fixed the reward-model random seed across partition draws, so
within-seed variation reflected the sampled partition rather than model
initialization. It used the same two generating setups as Direction 1 and 20
sampled partitions for FCM, PCM, and rough K-means.

\paragraph{Findings.}
Mean sampled estimate ranges were large, about 1.75--2.04 across methods and
generating setups. None of the sampled min--max or 5th--95th percentile ranges
covered the true value in any seed. Sampling generally worsened FCM and rough
K-means relMSE. PCM sampling improved over its degenerate hardened partition in
many draws, but remained noncompetitive. Sensitivity width alone did not reliably
predict OffCEM failure.

\subsection{Controlled FCM Overlap}

\paragraph{Purpose.}
This controlled experiment tested whether FCM membership uncertainty adds
information beyond policy-weighted DM when true action overlap is known.

\paragraph{Result path.}
\texttt{/Users/cispa/Documents/OffCEM/imprecise\_results/controlled\_fcm\_overlap/}.

\paragraph{Design.}
The generator created known two-cluster memberships and varied the fraction of
ambiguous actions. It used two reward modes:
\texttt{reward\_relevant}, where rewards mix latent cluster reward functions
using the known memberships, and \texttt{feature\_only}, where feature overlap is
present but rewards use only the primary hard cluster. The completed run covered
2 reward modes, 4 ambiguity levels, 10 seeds, and 1760 partition evaluations.

\paragraph{Findings.}
FCM did not recover the known memberships: ambiguity AUC was approximately
chance, estimated entropy was high even at zero true ambiguity, and membership
entropy had near-zero or negative correlation with the reward-mixture effect.
FCM-hard OffCEM error was nearly identical between reward-relevant and
feature-only modes, indicating that degradation came mainly from worsening FCM
partitions as features overlapped, not from successful detection of
reward-relevant ambiguity. Adding FCM uncertainty improved held-ambiguity error
calibration modestly but not ranking.

\subsection{Relaxed Local Correctness With Rough K-Means}

\paragraph{Purpose.}
This experiment evaluated an epsilon-relaxed local-correctness bias bound for
the population OffCEM functional while varying rough-partition ambiguity.

\paragraph{Result path.}
\texttt{/Users/cispa/Documents/OffCEM/imprecise\_results/relaxed\_local\_correctness/}.

\paragraph{Design.}
The run varied ambiguity fractions
$\{0,0.25,0.5,0.75\}$ and rough ratios
$\{1.05,1.15,1.25,1.5\}$ with 5 seeds. It evaluated the hardened rough
partition, sampled rough partitions, and the oracle primary partition. It
computed the fitted finite-sample OffCEM estimate, population OffCEM functional,
exact context-cluster residual oscillation epsilon, TV-weighted theoretical bias
bound, and coverage.

\paragraph{Findings.}
The exact epsilon-TV bound covered 100\% of all 560 partition evaluations. The
finite-sample estimate coverage was also 100\% in the grid. The bound was
conservative: about 2.1 times the actual absolute population bias, with interval
widths around 13--14 while the mean true policy value was 5.47. Rough K-means
detected increasing ambiguity, but sampled rough partitions added little to the
robust interval width. The result validates the bound mathematically in the
controlled setup, but does not yet provide a sharp practical interval.

\subsection{Domain-Informed and Outcome-Aware Clustering}

\paragraph{Purpose.}
This experiment tested whether OffCEM clusters should come from generic metadata,
external domain knowledge, auxiliary-data outcome profiles, or hybrid
representations.

\paragraph{Result path.}
\texttt{/Users/cispa/Documents/OffCEM/imprecise\_results/domain\_outcome\_clustering/}.

\paragraph{Design.}
The synthetic environment separated generic metadata, true domain classes,
observed noisy domain attributes, and latent outcome classes. Domain alignment
took values $\{0,0.5,1\}$ and domain-label noise took values
$\{0,0.25,0.5\}$, with 10 seeds per cell. Methods included feature-only K-means,
domain-informed K-means, oracle outcome-profile K-means, estimated
outcome-profile K-means from an independent auxiliary log, a hybrid
representation, shuffled-domain, and random controls.

\paragraph{Findings.}
Observable domain knowledge worked when it was reward-relevant. With clean
domain labels, domain-informed relMSE moved from 1.384 at zero domain alignment
to 0.454 at half alignment and 0.003 at full alignment. The benefit degraded
with label noise. Oracle outcome-profile clustering was the best clustering
overall and selected the lowest-error partition in 84/90 cells. Policy-weighted
DM selected the same argmin in 84/90 cells and had mean within-cell Spearman
correlation 0.77 with OffCEM squared error.

The estimated-outcome arm was inconclusive: it did not recover outcome structure
because the auxiliary reward model used generic action features that were
independent of domain and outcome labels. Thus the experiment validates clean
domain knowledge and oracle outcome geometry, but not the current estimated
outcome-profile implementation.

\section{Cross-Experiment Conclusions}

\begin{enumerate}
  \item Alignment between the estimation partition and reward residual geometry
  is the central condition for OffCEM success.
  \item The default synthetic benchmark is feature-arbitrary: the original
  generating partition is a stochastic softmax sample with near-zero alignment
  to K-means, and residual functions are keyed to cluster index.
  \item Local correctness degrades continuously under partition corruption;
  it is not an empirical all-or-nothing gate.
  \item The demeaned local-correctness metric predicts OffCEM error well within a
  fixed problem, but raw DM does not define a portable threshold across problem
  settings.
  \item Oracle-free recovery is possible when the clustering geometry is
  observable and reward-relevant. The feature-bucket benchmark and clean
  domain-informed experiment are the main positive examples.
  \item Current fuzzy/possibilistic/rough membership-uncertainty pilots do not
  improve OffCEM on the tested one-hot representations. Their strongest value so
  far is diagnostic, especially through policy-weighted DM.
  \item Epsilon-relaxed local correctness gives valid controlled bias bounds, but
  the current bounds are too wide to be practically sharp.
\end{enumerate}

\section{Completed-Result Inventory}

\begin{longtable}{p{0.28\linewidth}p{0.37\linewidth}p{0.27\linewidth}}
\toprule
Experiment & Completed result location & Status note \\
\midrule
\endhead
Clustering matched/mismatched &
\texttt{clustering\_results/} &
Completed, but older magnitudes should be regenerated due to later B1/B2 fixes. \\
Local correctness cluster-count sweep &
\texttt{local\_correctness\_results/\{matched,original,random,shuffled\_partition\}/} &
Authoritative local-correctness result source. \\
Local correctness temperature sweep &
\texttt{local\_correctness\_results/temp\_*} &
Completed for matched, kmeans, original, random. \\
Local correctness corruption sweep &
\texttt{local\_correctness\_results/corrupt\_p\_*} &
Completed continuum from matched to random. \\
Feature-coherent local correctness &
\texttt{local\_correctness\_results/feature\_coherent\_fb/} and
\texttt{feature\_coherent\_km/} &
Completed; confirms oracle-free recovery for feature-bucket geometry. \\
Join experiment, original generation &
\texttt{join\_results/cell\_*.json} &
Completed 36-cell grid. \\
Join experiment, feature-bucket generation &
\texttt{join\_results\_fc/cell\_*.json} &
Completed 36-cell grid. \\
Applicability Phase A smoke &
\texttt{applicability\_results/smoke/} &
Completed smoke only; not a final research conclusion. \\
Imprecise Direction 1 &
\texttt{imprecise\_results/direction1/} &
Completed pilot. \\
Imprecise Direction 6 corrected &
\texttt{imprecise\_results/direction6\_upd/} &
Completed corrected sensitivity pilot. \\
Controlled FCM overlap &
\texttt{imprecise\_results/controlled\_fcm\_overlap/} &
Completed. \\
Relaxed local correctness &
\texttt{imprecise\_results/relaxed\_local\_correctness/} &
Completed. \\
Domain/outcome clustering &
\texttt{imprecise\_results/domain\_outcome\_clustering/} &
Completed. \\
\bottomrule
\end{longtable}

\end{document}
